% =============================================================================
% Rozdział 6: Podsumowanie
% =============================================================================
\chapter{Podsumowanie}

\section{Podsumowanie zrealizowanych prac}

W ramach niniejszej pracy magisterskiej zrealizowano pełny cykl badawczy — od analizy literatury, przez projekt i~implementację systemu badawczego, do przeprowadzenia eksperymentów i~weryfikacji hipotezy na klastrze \textit{Google Kubernetes Engine} (GKE).

Przegląd literatury (rozdział 1) objął analizę istniejących badań nad autoskalowaniem w~\textit{Kubernetes}, w~wyniku której zidentyfikowano cztery luki badawcze: brak systematycznego porównania strategii HPA dla różnych profili obciążeniowych, ograniczoną analizę HPA Memory, brak badań łączących HPA z~modelem \textit{SaaS} oraz brak analizy zjawisk anomalnych w~autoskalowaniu chmurowym. Luki te bezpośrednio umotywowały kierunek dalszych prac.

Zaprojektowano system badawczy oparty na macierzy $2 \times 3$ (dwa profile obciążeniowe $\times$ trzy strategie HPA), składający się z~czterech mikroserwisów o~kontrolowanych profilach obciążeniowych, infrastruktury monitorującej opartej na \textit{kube-prometheus-stack} (Helm) oraz dedykowanej maszyny \textit{Compute Engine} jako generatora obciążenia (rozdział 3). Architektura została zoptymalizowana pod kątem izolacji zmiennych — każdy scenariusz testuje dokładnie jedną kombinację profilu i~strategii, przy pozostałych parametrach utrzymywanych na stałym poziomie.

Zaimplementowano kompletny system badawczy w~języku Python 3.12 z~frameworkiem \textit{FastAPI}, konteneryzowany za pomocą \textit{Docker} i~wdrażany na zarządzanym klastrze GKE (rozdział 4). Wszystkie manifesty \textit{Kubernetes} (Deploymenty, Serwisy typu \texttt{LoadBalancer}, konfiguracje HPA) oraz konfiguracje monitoringu są dostępne w~repozytorium, zapewniając pełną odtwarzalność badań. Opracowano skrypt \texttt{run-tests.sh} orkiestrujący wykonanie scenariuszy testowych z~automatycznym przełączaniem HPA, resetowaniem \textit{Deploymentów} i~cooldownem między scenariuszami, oraz zdefiniowano trzy skrypty \textit{k6} z~precyzyjnymi progami akceptacji i~profilami obciążenia dopasowanymi do charakterystyki aplikacji.

Wdrożono infrastrukturę monitorującą przez \textit{kube-prometheus-stack} (Helm), zapewniającą automatycznie skonfigurowany stos \textit{Prometheus}--\textit{Grafana} z~niestandardowymi metrykami przez \textit{prometheus-adapter}. Zdefiniowano procedurę eksportu pięciu zestawów danych CSV na scenariusz, które wraz z~wynikami \textit{k6} tworzą kompletny zbiór danych badawczych (rozdział 5).

Przeprowadzono badania eksperymentalne według zaplanowanej procedury (8 scenariuszy) na klastrze GKE, a~następnie dokonano analizy porównawczej strategii HPA z~wykorzystaniem siedmiu metryk badawczych. Szczególną uwagę poświęcono trzem zjawiskom anomalnym — \textit{HPA Lag}, \textit{Cold Start Latency} i~\textit{CPU Throttling Cascade} — dla których zaproponowano metody ilościowej charakterystyki na podstawie danych z~eksportów CSV z~\textit{Grafany}.

\section{Weryfikacja hipotezy badawczej}

Hipoteza badawcza postawiona w~rozdziale 1 brzmiała:

\begin{quote}
\textit{,,Nie istnieje jedna uniwersalna metryka HPA — skuteczność strategii skalowania zależy od charakterystyki obciążeniowej aplikacji (\textit{workload profile}). Błędny wybór metryki prowadzi do niedoskalowania lub przeskalowania.''}
\end{quote}

Wyniki eksperymentów \textbf{częściowo potwierdzają} postawioną hipotezę, jednak w~sposób bardziej złożony niż pierwotnie zakładano.

\textbf{HPA CPU okazał się jedyną skuteczną strategią} — wywołał skalowanie zarówno dla aplikacji CPU-bound (29 replik), jak i~I/O-bound (7 replik). Jest to wynik częściowo sprzeczny z~hipotezą, która przewidywała nieskuteczność HPA CPU dla I/O-bound. W~praktyce, przy 150 współbieżnych VUs, narzut \textit{Uvicorn} i~pętli zdarzeń \textit{asyncio} generował wystarczające obciążenie CPU, by przekroczyć próg 50\%. Oznacza to, że HPA CPU \textit{może} być skuteczny również dla aplikacji I/O-bound — pod warunkiem odpowiednio niskiego limitu CPU na Podzie i~wysokiej współbieżności. Jednak dla CPU-bound skuteczność ta była okupiona poważnym przeskalowaniem: 29 replik (190\% powyżej \texttt{maxReplicas}) przy jedynie 26\% wzroście przepustowości ($E = 0{,}01$).

\textbf{HPA Memory był nieskuteczny w~obu profilach} — zgodnie z~przewidywaniami. Zużycie pamięci w~zaimplementowanych aplikacjach jest statyczne i~nie koreluje z~liczbą żądań. Wynik ten potwierdza, że HPA Memory nie nadaje się jako jedyna metryka dla aplikacji o~stałym profilu pamięciowym, choć może być wartościowy dla systemów cache, kolejek komunikatów czy przetwarzania dużych plików.

\textbf{HPA Custom RPS nie wywołał skalowania} — co jest wynikiem \textit{pozornie} sprzecznym z~hipotezą, która przewidywała jego uniwersalną skuteczność. Przyczyną nie był jednak błędny wybór metryki, lecz niewłaściwa kalibracja progu: 100 RPS/Pod okazało się wartością zbyt wysoką dla obu serwisów (łączny RPS $\sim$65, co przy 1--2 replikach daje $\sim$32--65 RPS/Pod). Wynik ten prowadzi do istotnego uściślenia hipotezy:

\begin{quote}
\textit{Skuteczność strategii HPA zależy nie tylko od wyboru metryki, ale również od kalibracji progu. Błędny dobór progu może uniemożliwić skalowanie nawet przy prawidłowo wybranej metryce — i~odwrotnie, prawidłowo dobrany próg może uczynić skuteczną metrykę, która teoretycznie nie powinna działać dla danego profilu obciążeniowego.}
\end{quote}

Niezależnie od konkluzji dotyczącej hipotezy głównej, analiza zjawisk anomalnych dostarcza ilościowych danych na temat ograniczeń autoskalowania reaktywnego w~środowisku GKE — \textit{HPA Lag} ($\sim$60--75\,s) i~\textit{Cold Start Latency} ($\sim$15--25\,s) oznaczają, że system potrzebuje 1--2 minut na pełną reakcję na nagły wzrost obciążenia, co powinno być uwzględniane przy definiowaniu progów SLA.

\section{Ograniczenia pracy}

Należy wskazać następujące ograniczenia przeprowadzonych badań.

Aplikacje badawcze (\texttt{cpu-service}, \texttt{io-service}) symulują określone profile obciążeniowe, ale nie odzwierciedlają rzeczywistych systemów produkcyjnych, które często wykazują profile mieszane (\textit{CPU-bound} + \textit{I/O-bound}). Badanie ogranicza się do trzech metryk HPA — CPU, Memory i~Custom RPS — nie analizując innych strategii, takich jak skalowanie na podstawie liczby żądań w~toku (\textit{in-flight}), metryk zewnętrznych (np. długość kolejki \textit{RabbitMQ}) ani skalowania predykcyjnego opartego na modelach uczenia maszynowego. Wszystkie eksperymenty przeprowadzono w~pojedynczej strefie GCP (\texttt{europe-central2}), co oznacza, że nie badano wpływu opóźnień między strefami (\textit{cross-zone latency}) ani scenariuszy \textit{multi-region}. Progi HPA (CPU 50\%, Memory 70\%, RPS 100/Pod) zostały ustalone arbitralnie na podstawie wartości referencyjnych i~nie były przedmiotem optymalizacji — jak wykazały badania, kalibracja progu ma kluczowe znaczenie dla skuteczności strategii. Rozszerzenie \textit{SaaS Multi-Tenancy} (Faza VIII planu pracy) zostało zaprojektowane koncepcyjnie, ale nie zaimplementowane w~kodzie, co oznacza, że wnioski dotyczące zachowania systemu w~środowisku wielodostępnym pozostają hipotetyczne i~wymagają dalszych badań.

\section{Rekomendacje praktyczne}

Na podstawie przeprowadzonych badań sformułowano następujące rekomendacje dla praktyków projektujących polityki autoskalowania w~\textit{Kubernetes}.

Fundamentalną zasadą jest dopasowanie metryki HPA do profilu obciążeniowego aplikacji \textit{oraz} empiryczna kalibracja progu w~warunkach zbliżonych do produkcyjnych. Jak wykazały badania, nawet prawidłowo dobrana metryka (Custom RPS) nie gwarantuje skalowania przy źle skalibrowanym progu, a~metryka teoretycznie nieoptymalna (CPU) może okazać się skuteczna przy odpowiedniej konfiguracji limitów i~progów.

W~zakresie operacyjnym należy uwzględniać zjawiska anomalne przy definiowaniu SLA — \textit{HPA Lag} i~\textit{Cold Start Latency} oznaczają, że realistyczny czas reakcji systemu na nagły wzrost obciążenia wynosi 1--2 minuty, a~nie kilka sekund. \textit{CPU throttling} powinien być aktywnie monitorowany (metryka \texttt{container\_cpu\_cfs\_throttled\_seconds\_total}), ponieważ nie powoduje awarii, ale drastycznie wydłuża czasy odpowiedzi, działając jako ,,cichy zabójca'' wydajności. Dla testów obciążeniowych i~środowisk produkcyjnych rekomenduje się bezpośrednią ekspozycję serwisów przez \textit{GCP LoadBalancer} zamiast \texttt{port-forward}, co eliminuje wąskie gardło tunelowania przez \textit{API server} i~zapewnia realistyczne warunki sieciowe. Wreszcie, wdrożenie monitoringu przez \textit{kube-prometheus-stack} (Helm) redukuje złożoność konfiguracji i~zapewnia spójność wersji komponentów, co jest szczególnie istotne w~środowiskach produkcyjnych.

\section{Dalsze kierunki badań}

Niniejsza praca otwiera kilka obiecujących kierunków dalszych badań.

Najpilniejszym rozszerzeniem jest implementacja modelu \textit{SaaS Multi-Tenancy} (Faza VIII planu pracy) i~przeprowadzenie testów w~środowisku wielodostępnym na GKE. Zaprojektowana koncepcja przewiduje dwa poziomy izolacji: \textit{application-level} (tenanci rozróżniani nagłówkiem \texttt{X-Tenant-ID}, współdzielona infrastruktura) oraz \textit{namespace-per-tenant} (dedykowane przestrzenie nazw z~własnymi \textit{ResourceQuota}, \textit{NetworkPolicy} i~niezależnymi HPA). Środowisko GKE, z~automatycznym skalowaniem węzłów i~natywną integracją z~ekosystemem GCP, jest naturalną platformą do tych badań — eliminuje ograniczenia sprzętowe występujące na pojedynczym serwerze lokalnym i~umożliwia testowanie z~realistyczną liczbą tenantów. Kluczowymi pytaniami badawczymi są: czy HPA chroni ,,cichego'' tenanta przed degradacją wydajności spowodowaną przez ,,głośnego sąsiada'' (\textit{noisy neighbor}), czy zjawiska anomalne (\textit{HPA Lag}, \textit{Cold Start Latency}, \textit{CPU Throttling Cascade}) nasilają się w~środowisku wielodostępnym oraz jaki jest koszt infrastrukturalny izolacji \textit{namespace-per-tenant} w~porównaniu z~modelem \textit{application-level}.

Kolejnym, nowo zidentyfikowanym kierunkiem — wynikającym bezpośrednio z~przeprowadzonych badań — jest systematyczna analiza wpływu kalibracji progów HPA na skuteczność skalowania. Jak wykazały eksperymenty, próg 100 RPS/Pod dla HPA Custom RPS okazał się zbyt wysoki; interesującym pytaniem badawczym byłoby wyznaczenie optymalnych progów dla każdej kombinacji profilu obciążeniowego i~strategii HPA, potencjalnie z~wykorzystaniem technik automatycznego strojenia (np. \textit{Bayesian optimization}).

Dalszym kierunkiem jest rozszerzenie macierzy badawczej o~dodatkowe profile obciążeniowe i~strategie HPA. Włączenie \texttt{memory-service} jako trzeciego profilu (\textit{memory-bound}) umożliwiłoby weryfikację, czy HPA Memory staje się skuteczne dla aplikacji, których zużycie pamięci koreluje z~obciążeniem — pytanie, na które obecne badania nie dają odpowiedzi. Analiza metryk takich jak liczba żądań w~toku (\texttt{http\_requests\_inflight}) lub kombinacji wielu metryk w~jednym HPA mogłaby ujawnić strategie pośrednie między prostotą HPA CPU a~uniwersalnością HPA Custom RPS. Porównanie reaktywnego HPA z~podejściami predykcyjnymi (np. opartymi na LSTM), które przewidują przyszłe obciążenie i~skalują proaktywnie, mogłoby wskazać drogę do eliminacji problemu \textit{HPA Lag}.

W~warstwie infrastrukturalnej wartościowym uzupełnieniem byłaby analiza kosztowa różnych strategii HPA w~modelu \textit{pay-as-you-go} GCP — szczególnie w~kontekście \textit{multi-tenancy}, gdzie przeskalowanie (jak zaobserwowane 29 replik dla \texttt{cpu-cpu}) przekłada się bezpośrednio na wyższe koszty — oraz powtórzenie eksperymentów w~konfiguracji \textit{multi-zone} w~celu zbadania wpływu opóźnień między strefami na skuteczność HPA i~zjawiska anomalne. Docelowo, walidacja wniosków na rzeczywistych aplikacjach produkcyjnych (np. system rezerwacji, przetwarzanie obrazu w~pipeline ML) pozwoliłaby ocenić, na ile wyniki uzyskane na syntetycznych serwisach badawczych przenoszą się na systemy klasy enterprise.

\section{Podsumowanie końcowe}

Niniejsza praca magisterska stanowi kompletne studium empiryczne nad skutecznością strategii \textit{HorizontalPodAutoscaler} w~\textit{Kubernetes} dla różnych profili obciążeniowych, zrealizowane w~całości w~środowisku \textit{Google Cloud Platform}. Główny wkład pracy obejmuje: systematyczną analizę literatury identyfikującą cztery luki badawcze, w~tym nowo zidentyfikowaną lukę dotyczącą analizy zjawisk anomalnych w~autoskalowaniu; projekt macierzy badawczej $2 \times 3$ umożliwiającej izolowane porównanie strategii HPA; implementację czterech mikroserwisów o~kontrolowanych profilach obciążeniowych wraz z~pełną infrastrukturą monitorującą; automatyzację procesu badawczego — od wdrożenia po wykonanie testów i~eksport danych; oraz koncepcję rozszerzenia \textit{SaaS Multi-Tenancy} dla środowiska GKE, uwzględniającą specyfikę zarządzanego klastra i~mechanizmy izolacji \textit{Kubernetes}.

Przeprowadzone badania empiryczne dostarczają dowodów na to, że skuteczność strategii HPA zależy zarówno od wyboru metryki, \textit{jak i} od kalibracji jej progu. Głównym, nieprzewidzianym wynikiem jest to, że HPA CPU — a~nie HPA Custom RPS — okazał się strategią najbliższą ,,uniwersalnej'' w~warunkach testowych: skalował się zarówno dla CPU-bound (29 replik, eliminacja błędów z~37,70\% do 0,04\%), jak i~dla I/O-bound (7 replik, wzrost przepustowości o~59\%). Jednocześnie ujawniono poważny problem przeskalowania (29 replik przy limicie 10) oraz fundamentalne znaczenie kalibracji progu — HPA Custom RPS nie skalował się nie dlatego, że metryka jest nieskuteczna, ale dlatego, że próg 100 RPS/Pod był zbyt wysoki. Badania dostarczają również ilościowej charakterystyki zjawisk anomalnych (\textit{HPA Lag} $\sim$60--75\,s, \textit{Cold Start Latency} $\sim$15--25\,s), które mają istotne znaczenie praktyczne dla projektantów systemów \textit{SaaS} działających na \textit{Kubernetes} w~chmurze.

Otwarte, odtwarzalne środowisko badawcze dostępne w~repozytorium umożliwia niezależną weryfikację wyników i~stanowi fundament dla dalszych badań w~kierunku \textit{multi-tenancy}, optymalizacji progów i~skalowania predykcyjnego.
